\section{Worked example: reversible half-adder}

A half-adder adds two one-bit numbers without a carry input. Classical theory
defines two outputs:
\[
\mathrm{Sum}=A\oplus B,\qquad
\mathrm{Carry}=A\land B.
\]
The name \emph{half-adder} means it handles half of the full-adder interface:
there is no $C_{\rm in}$ input.

\subsection{Classical specification}

\begin{center}
\begin{tabular}{cc|cc|c}
\toprule
$A$&$B$&Carry&Sum&Binary interpretation\\
\midrule
0&0&0&0&$0+0=00_2$\\
0&1&0&1&$0+1=01_2$\\
1&0&0&1&$1+0=01_2$\\
1&1&1&0&$1+1=10_2$\\
\bottomrule
\end{tabular}
\end{center}
The Carry column is the high-order bit and Sum is the low-order bit. In the
last row, decimal $1+1=2$, written $10_2$ in binary.

\subsection{QPU protocol}

\textbf{Complete runnable protocol}
\begin{lstlisting}
PARAMS: A:state B:state

MAIN-PROCESS HalfAdder
CREATETOKEN -I Carry Sum
SET Carry 0p
SET Sum 0p

XOR -I A B -O Sum
AND -I A B -O Carry

MEASURE -I Carry
MEASURE -I Sum
RETURNVALS Carry Sum
\end{lstlisting}

\subsection{What every term is doing}

\begin{enumerate}
  \item \textbf{PARAMS} declares A and B as caller-controlled state inputs.
  In a basis-state truth-table test they receive \code{0p} or \code{1p}.
  \item \textbf{Carry and Sum are workspace targets.} CREATETOKEN reserves their
  names; SET prepares both in zero.
  \item \textbf{XOR uses A and B as controls/inputs.} It mutates Sum according
  to $s'=s\oplus A\oplus B$. Because $s=0$, the result is $A\oplus B$.
  \item \textbf{AND uses the same preserved controls.} It mutates Carry
  according to $c'=c\oplus(A\land B)$. Because $c=0$, Carry becomes the
  classical carry predicate.
  \item \textbf{MEASURE produces classical samples.} For basis inputs this
  example is deterministic. If A or B is in superposition, outputs may be
  entangled with the inputs and measurement samples a branch.
  \item \textbf{RETURNVALS defines interface order.} Carry is output column
  zero and Sum is output column one. It does not perform measurement itself.
\end{enumerate}

\begin{center}
\begin{tikzpicture}[scale=.82]
  \foreach \y/\name in {1.5/A,.5/B,-.5/Sum,-1.5/Carry} {
    \draw[thick] (0,\y)--(8,\y);
    \node[left] at (0,\y) {\name};
  }
  \fill[QpuNavy] (2,1.5) circle (.09);
  \fill[QpuNavy] (2,.5) circle (.09);
  \draw[thick] (2,1.5)--(2,-.5);
  \node[draw=QpuBlue,fill=QpuBlue!8,minimum width=.72cm,minimum height=.5cm,font=\scriptsize\bfseries] at (2,-.5) {XOR};
  \fill[QpuNavy] (4.2,1.5) circle (.09);
  \fill[QpuNavy] (4.2,.5) circle (.09);
  \draw[thick] (4.2,1.5)--(4.2,-1.5);
  \node[draw=QpuBlue,fill=QpuBlue!8,minimum width=.72cm,minimum height=.5cm,font=\scriptsize\bfseries] at (4.2,-1.5) {AND};
  \node[draw=QpuBlue,fill=QpuBlue!8,minimum width=.6cm,minimum height=.5cm,font=\bfseries] at (6.2,-.5) {M};
  \node[draw=QpuBlue,fill=QpuBlue!8,minimum width=.6cm,minimum height=.5cm,font=\bfseries] at (6.2,-1.5) {M};
  \draw[-{Latex[length=2mm]},QpuRed,thick] (6.55,-.5)--(7.8,-.5);
  \draw[-{Latex[length=2mm]},QpuRed,thick] (6.55,-1.5)--(7.8,-1.5);
\end{tikzpicture}
\end{center}

\begin{beginnerbox}[Which wires change?]
Only the two outputs. A and B are \emph{read} by both gates and come out
exactly as they went in. Sum and Carry are \emph{workspace}: they start at 0
and each receives one answer. Keeping the inputs is what makes the circuit
reversible, and it is why a quantum half-adder needs four wires where a
classical one draws two inputs and two outputs.
\end{beginnerbox}

\subsection{Build the Carry gate in the workbench}

The AND line of the protocol corresponds to these workbench choices. With A
on q0, B on q1, and Carry on q2:
\begin{center}
\WorkbenchControlMap{AND}
\end{center}
\begin{trybox}
Choose the \emph{Half adder} starter circuit. Set \textbf{q0 start} and
\textbf{q1 start} to each combination of \code{0p} and \code{1p}, choose
\textbf{Run all} each time, and check the measured Sum and Carry against the
table above.
\end{trybox}

\subsection{Reversible state equation}

Before measurement, the full basis-state action is
\[
\ket{A,B,s,c}\longmapsto
\ket{A,B,s\oplus A\oplus B,\ c\oplus(A\land B)}.
\]
The familiar half-adder is the special case $s=c=0$. Writing the complete
equation avoids a common mistake: XOR and AND do not overwrite arbitrary
targets. They reversibly combine results with the targets' previous states.

\subsection{Companion truth-table metadata}

\begin{lstlisting}
MAIN-PROCESS: HalfAdder
INPUTS: A B
OUTPUTS: Carry Sum
# A B Carry Sum
0 0p 0p 0p 0p
1 0p 1p 0p 1p
2 1p 0p 0p 1p
3 1p 1p 1p 0p
\end{lstlisting}
This file compares the measured QPU outputs with the classical half-adder
specification. Passing all four rows demonstrates basis-state equivalence; it
does not by itself characterize phases for superposition inputs.

